% Page 019

\seite{19}

Das Jahr 1891.\ob
Die Schülerzahl beträgt in diesem Jahre 43.\ Die landwirthschaftlichen\ob
Erträge und das Jahr sind mittelmäßig. Der Preis der meisten\ob
Erzeugnisse ließ zu wünschen übrig, besonders bei den\ob
Kartoffeln und dem Normandie-Roggen.

\pend
\abschnitt{Wallfahrt}
\pstart

Um die Wallfahrt zum St.\ Rochus auf dem heiligen Felsen alle\ob
alte Pfarreigeschichten. Es folgten 2 Prozessionen und\ob
zwar die erste am 23.\ Aug., die zweite am 18.\ Sept. je 3 Tage\ob
nach Trier.

Ein Gesuch, welches die hiesige Pfarrgemeinde um Einführung\ob
der Zulagsgehälfe die der königl. Regierung in Trier ein\ohb
reichte, wurde abschlägig beschieden.

Am 30.\ Mai 1891 schied der Lehrer Joh.\ Scholtes aus der Gemeinde,\ob
um die ihm übertragene Stelle zu Salmwald\unsicher{} in dem\ob
Kreise zu übernehmen. Sein Nachfolger wurde der Lehrer\ob
Joh.\ B. Halfen aus Dudeldorf, Kreis Bitburg. Dieser\ob
wurde in dem Seminar zu Wittlich vorgebildet und am\ob
2.\ November 1891 durch den Lokalschulinspector in feierlicher\ob
Weise eingeführt.

Im Jahre 1891 wurde der Gemeinde der Antrag gestellt, die\ob
Communicationsverbindung zwischen Hasselstein und Sefferweich\ob
herzustellen; da die Hauptkosten größtentheils auf Hasselstein fallen\ob
sollten; während die Gemeinde für die Grundabtretung\ob
und die weitere Instandhaltung der Straße zu Verfügung stellen\ob
eine geübte Verfolgung durch den Schultheiß\unsicher{} hierüber auf\ob
auf Antrag des Gemeinderates.
\pend
